\section{The fitted model and its spatial map}\label{sec:setting}
\subsection{Coordinates and ports}
Let $u\in\R^n$ be the FE displacement vector after the essential constraints
are imposed. Translations are in metres and rotations in radians. The
physical equations relate input forces $f\in\R^p$ to measured accelerations
$a\in\R^m$:
\begin{equation}
 M\ddot u+C\dot u+Ku=Bf,\qquad a=S\ddot u,
 \label{eq:physical}
\end{equation}
where $B\in\R^{n\times p}$ applies the loads and $S\in\R^{m\times n}$ samples
the responses. The true operators $M$, $C$ and $K$ are generally unknown. The
nominal FE prior supplies $M_0,K_0\succ0$. The reduced model retains $r\le n$
modal coordinates $q$ through the projection $u=Tq$ with $T\in\R^{n\times r}$
and obeys
\begin{equation}
 \ddot q+\Gamma\dot q+\Omega^2q=T^{\top}Bf,\qquad
 \Omega=\diag(\omega),\quad \omega_j>0,\quad\Gamma=\Gamma^{\top}\succ0.
 \label{eq:latent}
\end{equation}
The modal force $T^{\top}Bf$ and the physical load deliver the same power
$(Bf)^{\top}\dot u=(T^{\top}Bf)^{\top}\dot q$. A single matrix $T$ thus maps
displacements and its transpose maps forces. The coordinates $q$ are the
mass-normalized undamped modal coordinates of the fitted model. Its modal
mass $I_r$ and modal stiffness $\Omega^2$ are diagonal and leave the modal
damping matrix $\Gamma$ as the only coupling. These coordinates need not be
modal for the prior because neither $T^{\top}M_0T=I_r$ nor
$T^{\top}K_0T=\Omega^2$ is assumed. Translational rows of $T$ have units
kg$^{-1/2}$ and $q$ has units kg$^{1/2}$\,m. Both $\omega$ and $\Gamma$ have
units s$^{-1}$.

Let the full-row-rank matrix $A\in\R^{h\times n}$ sample the $h$ independent
directions of all registered ports (observers and force actuators). Fixed
selectors $R_{\rm o}\in\R^{m\times h}$ and $R_{\rm i}\in\R^{p\times h}$ pick
the response rows and the force rows of $A$:
\begin{equation}
 S=R_{\rm o}A,\qquad B^{\top}=R_{\rm i}A,\qquad W=AT.
 \label{eq:registration}
\end{equation}
The matrix $W\in\R^{h\times r}$ is the modal matrix $T$ restricted to the
ports. The selectors $R_{\rm o}$ and $R_{\rm i}$ pick disjoint rows of $A$ in
the examples because no port there both applies a force and measures a
response. But, in general, a port can do both, and then $R_{\rm o}$ and $R_{\rm i}$ share a row.

\subsection{Fit at the ports, complete away from them}
The model parameter is $\theta=(W,\omega,\Gamma)$: the port rows $W=AT$ of
the modal matrix, the undamped natural frequencies $\omega$ and the modal
damping $\Gamma$. Section~\ref{sec:identification} fits it to measured
accelerances.

Denote Laplace transforms by the argument $s$. With zero initial conditions
each time derivative becomes a factor $s$, so Eq.~\eqref{eq:latent} and
$a=S\ddot u=ST\ddot q$ become
\begin{equation}
 \Delta_\theta(s)q(s)=T^{\top}Bf(s),\qquad a(s)=s^2STq(s),\qquad
 \Delta_\theta(s)=s^2I_r+s\Gamma+\Omega^2.
 \label{eq:laplace}
\end{equation}
By Eq.~\eqref{eq:registration}, $T$ meets the measurement only at the ports:
\begin{equation}
 ST=R_{\rm o}AT=R_{\rm o}W,\qquad
 T^{\top}B=(AT)^{\top}R_{\rm i}^{\top}=W^{\top}R_{\rm i}^{\top}.
 \label{eq:ports-see}
\end{equation}
Substituting and solving for $q$ wherever $\Delta_\theta(s)$ is invertible,
\begin{equation}
 q(s)=\Delta_\theta(s)^{-1}W^{\top}R_{\rm i}^{\top}f(s),\qquad
 a(s)=H_\theta(s)f(s),\qquad
 H_\theta(s)=s^2R_{\rm o}W\Delta_\theta(s)^{-1}W^{\top}R_{\rm i}^{\top}.
 \label{eq:model}
\end{equation}
At $s=\mathrm i2\pi\nu$ with frequency $\nu$ in Hz, $H_\theta$ is the
accelerance FRF. It depends on $T$ only through $W$. So every $T$ with
$AT=W$ reproduces it. The fitted $\theta$ thus gives $q$ from the measured
forces. Only $T$ is left open, and it distributes $q$ over the mesh as
$u=Tq$. Figure~\ref{fig:spaces} shows the two routes from $f$ to $a$.

\begin{figure}[tb]\centering
% Figure 1: the accelerance composed through the FE space (top) and the port space (bottom).
% Blue maps are fitted from the FRFs, orange maps are completed with the nominal
% model, black maps are the fixed registration. Plain TikZ coordinates keep every
% label clear of every arrow.
\begin{tikzpicture}[x=1cm,y=1cm,>=stealth,
 sp/.style={align=center,inner sep=2pt,font=\small},
 lab/.style={font=\footnotesize,inner sep=1.5pt},
 fit/.style={->,thick,draw=blue!70!black},
 cmp/.style={->,thick,draw=orange!80!black},
 reg/.style={->,draw=black!75}]
 \newcommand{\spname}[2]{$#1$\\[-2pt]{\scriptsize\color{black!60}#2}}
 % nodes: forces -> FE load / port forces -> modal force -> modal displacement
 %        -> FE displacement / port displacement -> accelerations
 \node[sp] (f)  at (0,0)       {\spname{f\in\C^{p}}{measured forces}};
 \node[sp] (Bf) at (2.7,1.45) {\spname{Bf\in\C^{n}}{FE load}};
 \node[sp] (Rf) at (2.7,-1.45){\spname{R_{\rm i}^{\top}f\in\C^{h}}{port forces}};
 \node[sp] (g)  at (5.35,0)    {\spname{g\in\C^{r}}{modal force}};
 \node[sp] (q)  at (8.55,0)    {\spname{q\in\C^{r}}{modal displacement}};
 \node[sp] (u)  at (11.3,1.45) {\spname{u=T_\theta q\in\C^{n}}{FE displacement}};
 \node[sp] (y)  at (11.3,-1.45){\spname{y=Wq\in\C^{h}}{port displacement}};
 \node[sp] (a)  at (14.0,0)    {\spname{a/s^{2}\in\C^{m}}{accelerations}};
 % registration (black)
 \draw[reg] (f) -- node[lab,above left]{$B$} (Bf);
 \draw[reg] (f) -- node[lab,below left]{$R_{\rm i}^{\top}$} (Rf);
 \draw[reg] (Rf) -- node[lab,right]{$A^{\top}$} (Bf);
 \draw[reg] (u) -- node[lab,left]{$A$} (y);
 \draw[reg] (u) -- node[lab,above right]{$S$} (a);
 \draw[reg] (y) -- node[lab,below right]{$R_{\rm o}$} (a);
 % completed (orange) and fitted (blue)
 \draw[cmp] (Bf) -- node[lab,above right]{$T_\theta^{\top}$} (g);
 \draw[fit] (Rf) -- node[lab,below right]{$W^{\top}$} (g);
 \draw[fit] (g) -- node[lab,above]{$\Delta_\theta(s)^{-1}$} (q);
 \draw[cmp] (q) -- node[lab,above left]{$T_\theta$} (u);
 \draw[fit] (q) -- node[lab,below left]{$W$} (y);
 % Completion of the orange maps: three conditions on T_theta.
 \node[draw=orange!80!black,rounded corners=2pt,fill=orange!4,inner sep=3pt]
  (cmp) at (7.0,2.75)
 {\small\begin{tabular}{@{}c@{\hspace{2.2em}}c@{\hspace{2.2em}}c@{}}
  $AT_\theta=W$ & $\rank T_\theta=r$ & $(M_0^{-1}K_0)\,T_\theta\approx T_\theta\Omega^2$\\[-1pt]
  {\scriptsize\color{black!60}port rows kept} &
  {\scriptsize\color{black!60}every modal coordinate kept} &
  {\scriptsize\color{black!60}nominal undamped dynamics}
 \end{tabular}};
 \node[font=\scriptsize,text=orange!80!black,anchor=south west,inner sep=1pt]
  at ([xshift=2pt]cmp.north west) {completion of the orange maps};
 % legend
 \node[font=\scriptsize] at (7.0,-2.45)
 {\tikz[baseline=-0.6ex]\draw[fit](0,0)--(0.5,0); fitted from the FRFs\quad
  \tikz[baseline=-0.6ex]\draw[cmp](0,0)--(0.5,0); completed with $(M_0,K_0)$\quad
  \tikz[baseline=-0.6ex]\draw[reg](0,0)--(0.5,0); registration: $B=A^{\top}R_{\rm i}^{\top}$, $S=R_{\rm o}A$};
\end{tikzpicture}

\caption{The accelerance $H_\theta(s)$ composed through the FE space (upper
route) and through the port space (lower route). Both routes share the modal
force $g=T_\theta^{\top}Bf=W^{\top}R_{\rm i}^{\top}f$ and the modal
displacement $q$, and every cell commutes, since $AT_\theta=W$,
$B=A^{\top}R_{\rm i}^{\top}$ and $S=R_{\rm o}A$. The FRFs determine the lower
route and with it $q$. Completion supplies the upper route through
$T_\theta$. It keeps the port rows and all $r$ modal coordinates. It minimizes the
nominal $F=ma$ residual of $u=T_\theta q$ under the measured forces.}
\label{fig:spaces}
\end{figure}

\subsection{Pole--residue form of the fit}
The experiment measures the accelerance matrix
$\widehat H_\ell\in\C^{m\times p}$ at $N$ frequency lines
$s_\ell=\mathrm i2\pi\nu_\ell$. Each entry is a complex transfer from one
force to one response. The fit uses all of them. It regresses the model
$H_\theta(s_\ell)$ onto $\widehat H_\ell$ at every line.

The regression parameters are the number of modes $r$ and the model
parameter $\theta$. The order $r$ is set first. The laboratory example
selects it by validation (Section~\ref{sec:lab}). The least squares then
fits $\theta$ for that $r$.

The poles $\lambda_k$ of the model are the $2r$ roots of
$\det\Delta_\theta(\lambda)=0$. They depend on $\omega$ and $\Gamma$ only.
We assume every root is simple: no two roots coincide. Each simple root then
has one eigenvector direction. We take $\phi_k$ with
$\Delta_\theta(\lambda_k)\phi_k=0$ and normalize it by
$\phi_k^{\top}(2\lambda_kI_r+\Gamma)\phi_k=1$. Then $\Delta_\theta^{-1}$ is a
sum of rank-one terms over the poles \citep{lancaster1966,tisseur2001}.
Substituting it into Eq.~\eqref{eq:model} gives the form that the fit
evaluates:
\begin{equation}
 \begin{aligned}
 \Delta_\theta(s)^{-1}&=\sum_{k=1}^{2r}\frac{\phi_k\phi_k^{\top}}{s-\lambda_k},\\
 H_\theta(s)&=R_{\rm o}WW^{\top}R_{\rm i}^{\top}
 +\sum_{k=1}^{2r}\frac{\lambda_k^2(R_{\rm o}W\phi_k)(R_{\rm i}W\phi_k)^{\top}}{s-\lambda_k}.
 \end{aligned}
 \label{eq:poleresidue}
\end{equation}
The constant term is the acceleration feedthrough. All $mp$ transfers are
built from the same $2r$ poles. At pole $\lambda_k$ the residue is the
product of one column $\lambda_k^2R_{\rm o}W\phi_k$ over the responses and
one row $(R_{\rm i}W\phi_k)^{\top}$ over the forces. So every measured
transfer at every line constrains the same parameters $\omega$, $\Gamma$ and
$W$.

\subsection{Two steps}
The unknowns are $\theta$ and the modal matrix $T$. They are linked by
$AT=W$. The FRFs enter through a misfit $J(\theta)$ defined in
Section~\ref{sec:identification}. The nominal model $(M_0,K_0)$ enters
through a criterion $J_F(\theta,T)$ defined in Section~\ref{sec:completion}.
A joint estimate would solve
\begin{equation}
 \min_{\theta,\,T}\ J(\theta)+\tau J_F(\theta,T)
 \quad\text{subject to}\quad AT=W,\quad\rank T=r,\qquad\tau>0.
 \label{eq:joint}
\end{equation}
The rank condition keeps all $r$ modal coordinates on the mesh. By
Eqs.~\eqref{eq:ports-see} and \eqref{eq:model} the FRFs see $T$ only
through $W$. So $J$ does not depend on the rest of $T$. For fixed $\theta$
the best $T$ minimizes $J_F$ alone. It is the same for every $\tau$. The
weight $\tau$ acts on $\theta$ only. It lets the nominal model pull $\theta$
away from the measured FRFs. We give the data priority and let
$\tau\to0^+$. The joint problem then splits into two steps:
\begin{equation}
 \text{(i)}\quad\min_\theta J(\theta),\qquad
 \text{(ii)}\quad\min_{\substack{AT=W\\\rank T=r}}J_F(\theta,T)\ \text{ at the fitted }\theta.
 \label{eq:two-steps}
\end{equation}
Step~(i) uses the ports alone (Section~\ref{sec:identification}). Step~(ii)
uses the FE operators (Section~\ref{sec:completion}). It keeps $AT=W$, so it
never changes the fitted transfers.
